\documentclass[11pt]{article}
\usepackage[T1]{fontenc}
\usepackage[utf8]{inputenc}
\usepackage{lmodern,microtype}
\usepackage[margin=1in]{geometry}
\usepackage{amsmath,amssymb,amsthm,mathtools}
\usepackage{booktabs,array}
\usepackage{xcolor,listings,fancyhdr,xurl}
\usepackage[colorlinks=true,linkcolor=blue!45!black,citecolor=blue!45!black,urlcolor=blue!50!black]{hyperref}
\hypersetup{pdftitle={Desargues and Pascal in intuitionistic primitive recursive arithmetic},
pdfauthor={Prepared for Vladimir Reshetnikov},
pdfsubject={Complete constructive coordinate proofs, iPRA interpretation, and Lean formalization plan},
pdfkeywords={Desargues, Pascal, projective geometry, iPRA, constructive mathematics, Lean}}
\definecolor{ink}{RGB}{24,47,71}
\definecolor{codebg}{RGB}{247,248,250}
\pagestyle{fancy}
\fancyhf{}
\fancyhead[L]{\small Desargues and Pascal in iPRA}
\fancyhead[R]{\small Constructive proof package}
\fancyfoot[C]{\thepage}
\setlength{\headheight}{14pt}
\setlength{\emergencystretch}{3em}
\lstset{basicstyle=\ttfamily\small,backgroundcolor=\color{codebg},frame=single,
breaklines=true,columns=fullflexible,keepspaces=true,showstringspaces=false}
\newtheorem{theorem}{Theorem}[section]
\newtheorem{lemma}[theorem]{Lemma}
\newtheorem{proposition}[theorem]{Proposition}
\newtheorem{corollary}[theorem]{Corollary}
\theoremstyle{definition}
\newtheorem{definition}[theorem]{Definition}
\theoremstyle{remark}
\newtheorem{remark}[theorem]{Remark}
\newcommand{\N}{\mathbb N}
\newcommand{\Z}{\mathbb Z}
\newcommand{\Q}{\mathbb Q}
\newcommand{\R}{\mathbb R}
\newcommand{\ipra}{\mathsf{iPRA}}
\newcommand{\nz}{\operatorname{NZ}}
\newcommand{\br}[3]{[#1\,#2\,#3]}
\newcommand{\axis}{\operatorname{axis}}
\newcommand{\Zero}{\operatorname{Zero}}
\newcommand{\Apart}{\operatorname{Apart}}
\title{\color{ink}\textbf{Desargues and Pascal in intuitionistic\\
primitive recursive arithmetic}\\[0.5em]
\large Exact polynomial certificates, constructive geometry,\\
and a Lean formalization plan}
\author{Prepared for Vladimir Reshetnikov}
\date{15 November 2026}
\begin{document}
\maketitle
\begin{abstract}
We give complete algebraic proofs of Desargues's theorem, including its
converse, and Pascal's theorem, and explain their derivation in first-order
intuitionistic primitive recursive arithmetic. The literal arithmetic
theorems concern the rational projective plane, represented by nonzero
integer triples. Their conclusions include explicit primitive recursive
functions producing a common line or center. The principal certificates
are identities of integer polynomials: a factorization of the Desargues
determinant and an equality of the Pascal determinant with a six-by-six
quadratic evaluation determinant. The Pascal proof covers every nonzero
homogeneous quadratic, including singular conics, without a
general-position argument or a limiting argument. A further conditional
first-order schema covers primitive-recursive real coordinate names with
explicit apartness assumptions. The archive contains exact executable
polynomial checks and concrete Lean source candidates. The new Lean
sources have not been compiled, and a native iPRA derivation object has
not been emitted; the remaining formalization steps are specified
separately from the completed written proofs.
\end{abstract}
\clearpage
{\small\tableofcontents}
\newpage
\section{The precise result and its scope}

The usual Desargues theorem says that two nondegenerate triangles are
perspective from a point if and only if they are perspective from a line.
Pascal says that the three intersections of opposite sides of a hexagon
inscribed in a conic are collinear. Homogeneous coordinates include points
at infinity, so parallel affine lines require no exceptional case.
The statements in the question provide the geometric motivation
\cite{desargues,pascal}; all proofs below are supplied explicitly.

Three levels of statement must be distinguished in a first-order arithmetic
treatment.
\begin{enumerate}
\item \textbf{Internal arithmetic theorems.} For arbitrary integer-coded
coordinates and coefficients, iPRA proves the incidence implications and
the correctness of the displayed line and center algorithms. This is
the entire rational projective plane $\mathbb P^2(\Q)$, rather than a
finite numerical sample.
\item \textbf{Universal algebraic certificates.} The two main identities
hold in every commutative ring. Their geometric consequences hold over
commutative fields with the requisite nonzero, or constructive apartness,
witnesses. The proofs use no order or analytic structure.
\item \textbf{A real-coordinate arithmetic schema.}
Section~\ref{sec:real} states the same conclusions for each fixed finite
family of primitive-recursive rational approximation functions satisfying
explicit arithmetic regularity and apartness premises. This is an iPRA
schema, not quantification over a completed set of all real numbers.
\end{enumerate}
No assertion is made that Desargues follows from the incidence axioms of
an arbitrary abstract projective plane: coordinates are part of the
interpretation. Nor is an ambient Lean theorem about a ring automatically
identified with a derivation in the object theory iPRA.

Here are the principal statements; their notation and witness algorithms
are developed below. A triple is \emph{nonzero} when at least one of its
integer components is nonzero.

\begin{theorem}[Internal Desargues, in both directions]\label{thm:des-summary}
Let $a,b,c,A,B,C$ be integer triples with
$\br a b c\ne0$ and $\br A B C\ne0$. Put
\[
 j_1=a\times A,\quad j_2=b\times B,\quad j_3=c\times C,
\]
\[
 X=(a\times b)\times(A\times B),\quad
 Y=(b\times c)\times(B\times C),\quad
 Z=(c\times a)\times(C\times A).
\]
Suppose all $j_i$ and $X,Y,Z$ are nonzero. Then iPRA proves
\begin{align*}
 &\exists o\,[\nz(o)\land o\cdot j_1=o\cdot j_2=o\cdot j_3=0]\\
 &\qquad\longleftrightarrow
 \exists h\,[\nz(h)\land h\cdot X=h\cdot Y=h\cdot Z=0].
\end{align*}
In the forward direction choose $h=\axis(X,Y,Z)$; in the reverse
direction choose $o=\axis(j_1,j_2,j_3)$. The forward implication does not
need the two triangle determinant assumptions.
\end{theorem}

\begin{theorem}[Internal Pascal]\label{thm:pas-summary}
Let $P_1,\ldots,P_6$ be nonzero integer triples and let
\[
 Q(x,y,z)=q_1x^2+q_2y^2+q_3z^2+q_4xy+q_5xz+q_6yz
\]
have an integer coefficient vector $q\ne0$. Suppose $Q(P_i)=0$ for every
$i$. Define
\[
 X=(P_1\times P_2)\times(P_4\times P_5),\quad
 Y=(P_2\times P_3)\times(P_5\times P_6),
\]
\[
 Z=(P_3\times P_4)\times(P_6\times P_1).
\]
Then iPRA proves $\br X Y Z=0$. If $X,Y,Z$ are nonzero, it additionally
proves that $h=\axis(X,Y,Z)$ is nonzero and satisfies
$h\cdot X=h\cdot Y=h\cdot Z=0$.
\end{theorem}

The nonzero double cross products ensure that the side lines and their
specified intersections really exist. The determinant conclusion remains
true even when such objects degenerate, but the zero vector is never
treated as a projective point. Coincident intersections are allowed; a
common line is then constructed, although it need not be unique.
Replacing a repeated vertex by a tangent is a different construction and
is not silently included in the cross-product convention.


\section{The arithmetic foundation}\label{sec:arith}

\subsection{The exact version of iPRA}
We use first-order intuitionistic logic with equality, a symbol for every
primitive-recursive definition, its defining equations, successor
injectivity, $Sx\ne0$, and induction for quantifier-free formulas with
parameters. This is the system explicitly defined by
\texttt{PR}, \texttt{PTm}, and \texttt{PRAAx} in ProveIt's
\texttt{ConciseIPRA/PR.lean} \cite{repo-pr}.
Quantifier-free induction is not induction for arbitrary formulas about
geometric objects or real names.

Natural-number equality is decidable in this theory. Define its Boolean
test by primitive recursion, first using the zero test and predecessor,
and then comparing successors. The defining equations and
quantifier-free induction prove the test correct. Consequently every
finite Boolean combination of the integer equalities below permits
intuitionistically valid case analysis. This uses a decision procedure
for the specified predicate, rather than excluded middle for all formulas.

\subsection{Signed integers without a quotient-set assumption}
An integer is represented by a pair $(p,n)$ of natural numbers, read as
$p-n$. There is no need to form a set-theoretic quotient internally. Set
\begin{align*}
 (p,n)\simeq(p',n')&\quad\Longleftrightarrow\quad p+n'=p'+n,\\
 (p,n)+(r,s)&=(p+r,n+s),\\
 -(p,n)&=(n,p),\\
 (p,n)(r,s)&=(pr+ns,ps+nr).
\end{align*}
Zero and one are $(0,0)$ and $(1,0)$. These operations and the equality
test are primitive recursive and respect $\simeq$. Subsequent integer
equalities abbreviate this numerical equality relation. A tuple may be
left as finitely many variables or packed by a primitive-recursive
pairing function.

For completeness, the elementary arithmetic behind this interpretation
is as follows. Define $x+0=x$, $x+S y=S(x+y)$, $x\cdot0=0$, and
$x\cdot S y=x\cdot y+x$. Quantifier-free induction proves
$0+x=x$, $Sx+y=S(x+y)$, and associativity of addition. The first two
equations give commutativity of addition by induction on its second
argument. Induction on the repeated summand proves distributivity;
using distributivity, induction proves associativity and commutativity
of multiplication. Each induction formula is an equation with the other
quantities as parameters. Cancellation of a common additive summand
follows by induction on that summand and successor injectivity.
Expanding the pair operations now proves all commutative-ring laws:
both sides reduce to equal sums of natural monomials. The same expansions
show that replacing an input by an equivalent pair preserves the output.

We also need that these interpreted integers have no zero divisors.
Truncated subtraction $x\mathbin{\dot-}y$ is primitive recursive; its
recursion on $y$ uses predecessor. Its elementary defining identities
show that $(p,n)$ is equivalent to
\[
 (p\mathbin{\dot-}n,\ n\mathbin{\dot-}p),
\]
and that at least one of these two components is zero. One direct proof
uses the primitive-recursive comparison: strip successors from both
inputs until a zero is reached; the process can be bounded by $p+n$.
Induction on this explicit counter, with equations for the two registers
and the Boolean comparison result, is quantifier free. It proves the
usual alternatives $p=n+d$ or $n=p+d$, where the appropriate $d$ is the
displayed truncated difference. In either alternative the normalization
and its equality to $(p,n)$ follow by additive cancellation.

Every nonzero integer therefore has either the form $(S r,0)$ or
$(0,S r)$. The product of two such integers has a nonzero component,
since
\[
 S r\cdot S s=S\bigl(r+s+rs\bigr).
\]
This is unequal to zero by successor separation. Splitting on the
decidable zero test thus proves
\begin{equation}\label{eq:cancel}
 uv=0\ \land\ u\ne0\quad\Longrightarrow\quad v=0.
\end{equation}
All normalizations and comparisons here have primitive-recursive
definitions. None searches without a bound, and all the indicated
induction formulas are quantifier free.

\subsection{Why integer triples cover rational geometry}
A rational number can be presented by a signed numerator and a positive
natural denominator. Multiplying a triple by the product of its three
denominators produces an integer triple. Applying the same operation
to the six conic coefficients produces an integer quadratic equation.
No greatest common divisor or canonical representative is needed.
Multiplication by these nonzero denominators preserves incidence and
the zero or nonzero status of every determinant and cross product in
the statements. These operations are primitive recursive.

Thus rational points, rational lines, and rational homogeneous conics
can all be treated with the signed-pair arithmetic above. The later
real-name schema explains a separate interpretation for coordinates
or coefficients given by effective real approximations.


\section{Projective geometry as finite algebra}\label{sec:geometry}

\subsection{Conventions and elementary identities}
Vectors are columns, and coordinates are numbered $0,1,2$ in algorithms.
Set
\[
 u\cdot v=u_0v_0+u_1v_1+u_2v_2,
\]
\[
 u\times v=(u_1v_2-u_2v_1,\ u_2v_0-u_0v_2,\ u_0v_1-u_1v_0),
 \qquad \br u v w=u\cdot(v\times w)=\det(u,v,w).
\]
Two nonzero triples represent the same projective point when their cross
product is zero; lines have the same representation convention. Over
$\Q$ this agrees with equality up to nonzero scalar multiplication:
choose a nonzero coordinate of the first triple, and the cross-product
equations give the scalar ratio in every coordinate. The relation and
all incidences are invariant under changing representatives.

A point $p$ lies on a line $\ell$ precisely when $\ell\cdot p=0$.
The line joining distinct points $p,q$ is $p\times q$; the point where
distinct lines $\ell,m$ meet is $\ell\times m$. Orthogonality verifies
both incidences. For uniqueness, if $\ell\cdot p=\ell\cdot q=0$, then
\[
 \ell\times(p\times q)=p(\ell\cdot q)-q(\ell\cdot p)=0.
\]
Any nonzero such $\ell$ therefore represents the same line as $p\times q$.
The same identity proves the dual assertion for intersections.

We use determinant multilinearity, sign change on swapping columns,
zero with a repeated column, multiplicativity, and cofactor expansion.
There is no unbounded linear-algebra theorem hidden here. For the fixed
sizes $3,4,5,6$, define
\begin{equation}\label{eq:detperm}
 \det M=\sum_{\sigma\in S_n}\operatorname{sgn}(\sigma)
                   \prod_{i=1}^n M_{i,\sigma(i)}.
\end{equation}
The sum has at most $720$ terms. Multilinearity follows by distributivity.
Swapping columns pairs permutation terms with opposite signs; with equal
columns these terms cancel. Grouping terms by the selected entry in a
row or column gives cofactor expansion. Expanding $\det(UV)$, terms with
repeated selected indices cancel by the same pairing; the remaining
selections are permutations, giving $\det(UV)=\det U\det V$.
These are finite polynomial calculations over any commutative ring.

Two useful consequences, also checked directly from the coordinates, are
\begin{equation}\label{eq:doublecross}
 (u\times v)\times(r\times s)=\br u v s\,r-\br u v r\,s
\end{equation}
and the adjugate identity
\begin{equation}\label{eq:adj3}
 \br u v w\,o
 =(o\cdot u)(v\times w)+(o\cdot v)(w\times u)
                         +(o\cdot w)(u\times v).
\end{equation}
For the latter, the coefficient of each coordinate of $o$ is a cofactor
expansion; off-diagonal coefficients are determinants with a repeated
row. Thus its verification is the fixed-size argument just given.

\begin{lemma}[Incidence gives determinant zero]\label{lem:inc-det}
If $o$ is a nonzero integer triple and
$o\cdot u=o\cdot v=o\cdot w=0$, then $\br u v w=0$.
\end{lemma}
\begin{proof}
The right side of \eqref{eq:adj3} is zero. Select a nonzero coordinate
$o_i$ by testing the three coordinates. Its left coordinate is
$\br u v w\,o_i=0$. Apply \eqref{eq:cancel}.
\end{proof}

\subsection{An actual common line, not just a determinant test}
Define the total function
\begin{equation}\label{eq:axis}
\axis(u,v,w)=
\begin{cases}
 u\times v,&\nz(u\times v),\\
 u\times w,&\neg\nz(u\times v)\text{ and }\nz(u\times w),\\
 (-u_1,u_0,0),&\text{both crosses zero and }(u_0\ne0\text{ or }u_1\ne0),\\
 (u_2,0,0),&\text{otherwise}.
\end{cases}
\end{equation}
Here $\nz(u)$ abbreviates $u_0\ne0\lor u_1\ne0\lor u_2\ne0$.
Each branch uses only integer arithmetic and a fixed number of equality
tests, so this function is primitive recursive.

\begin{lemma}[Constructive common-kernel lemma]\label{lem:axis}
If $u\ne0$ and $\br u v w=0$, then $h=\axis(u,v,w)$ is nonzero and
$h\cdot u=h\cdot v=h\cdot w=0$.
\end{lemma}
\begin{proof}
In the first branch, orthogonality gives the first two equations and
$h\cdot w=\br u v w=0$. In the second, the analogous remaining equation
is $h\cdot v=-\br u v w=0$. Nonzeroness is the respective branch test.
In the third, $h\cdot u=0$, while
$h\cdot v=(u\times v)_2=0$ and
$h\cdot w=(u\times w)_2=0$; the test on $u_0,u_1$ proves $h\ne0$.
In the last branch, $u_0=u_1=0$ and $u_2\ne0$, because $u\ne0$.
Then $h\cdot u=0$ and
$h\cdot v=u_2v_0=(u\times v)_1=0$, with the same argument for $w$.
The component $u_2$ proves $h\ne0$.
\end{proof}
Applied to points, this constructs a common line; applied to lines, it
constructs a common point. It handles repeated nonzero inputs. Thus a
vanishing determinant is not being used as a definition that conceals the
existence assertion in either geometric theorem.


\section{Desargues: the identity and both implications}\label{sec:des}

\begin{lemma}[The universal Desargues identity]\label{lem:des-id}
With the notation of Theorem~\ref{thm:des-summary}, over every
commutative ring,
\begin{equation}\label{eq:des-id}
 \boxed{\br X Y Z=\br a b c\,\br A B C\,
              \br{a\times A}{b\times B}{c\times C}.}
\end{equation}
\end{lemma}
\begin{proof}
By \eqref{eq:doublecross},
\begin{align*}
 X&=\br a b B\,A-\br a b A\,B,\\
 Y&=\br b c C\,B-\br b c B\,C,\\
 Z&=\br c a A\,C-\br c a C\,A.
\end{align*}
Multilinearity gives eight terms. Six have repeated columns, leaving
\[
 \br X Y Z=E\br A B C,
\]
\[
 E=\br a b B\br b c C\br c a A
       -\br a b A\br b c B\br c a C.
\]
Let $T=(a\ b\ c)$ and $J=(a\times A\ \ b\times B\ \ c\times C)$.
Their matrix product is
\[
 T^{\mathsf t}J=
 \begin{pmatrix}
 0&\br a b B&-\br c a C\\
 -\br a b A&0&\br b c C\\
 \br c a A&-\br b c B&0
 \end{pmatrix}.
\]
Its diagonal is zero; the two surviving permutation terms in its
determinant are exactly those defining $E$. Multiplicativity gives
$E=\det T\det J$, proving \eqref{eq:des-id}. No cancellation or
nondegeneracy assumption entered this identity.
\end{proof}

\begin{proof}[Proof of Theorem~\ref{thm:des-summary}]
Suppose a nonzero center $o$ is given. Its three incidences and
Lemma~\ref{lem:inc-det} give $\br{j_1}{j_2}{j_3}=0$.
Equation~\eqref{eq:des-id} gives $\br X Y Z=0$.
Since $X\ne0$, Lemma~\ref{lem:axis} constructs the nonzero
axis $h=\axis(X,Y,Z)$. This proves the forward implication.

Conversely suppose a nonzero line $h$ through $X,Y,Z$ is given.
Lemma~\ref{lem:inc-det}, with $h$ in place of $o$, gives
$\br X Y Z=0$. Both $\br a b c$ and $\br A B C$ are nonzero,
so two applications of integer cancellation in \eqref{eq:des-id}
give $\br{j_1}{j_2}{j_3}=0$. Since $j_1\ne0$, the same algorithm
applied dually produces $o=\axis(j_1,j_2,j_3)$, a nonzero common center.
Every existential introduction uses this displayed term; every
existential elimination merely opens the supplied witness.
\end{proof}

\subsection{An exact example}
Take
\[
 a=(1,0,1),\quad b=(0,1,1),\quad c=(-1,-1,1),
\]
\[
 A=(1,0,2),\quad B=(0,1,3),\quad C=(-1,-1,4).
\]
The triangle determinants are $3$ and $9$. The joining lines are
$(0,-1,0)$, $(2,0,0)$, and $(-3,3,0)$, all through $(0,0,1)$.
The side intersections computed by the cross products are
\[
 X=(2,-1,1),\quad Y=(2,5,1),\quad Z=(-4,-1,-2).
\]
All lie on $(1,0,-2)$, or $x=2$ in the chart $z=1$.
The algorithm returns $X\times Y=(-6,0,12)$, representing that line.

\subsection{Why the converse needs nondegenerate triangles}
For
\[
 a=(1,0,0),\ b=(0,1,0),\ c=(1,1,0),\qquad
 A=(0,0,1),\ B=(1,0,1),\ C=(0,1,1),
\]
the side intersections are $(-1,0,0),(-1,1,0),(0,-1,0)$,
so are collinear. All joins and meets are nonzero. But the
joining-line determinant is $2$, so Lemma~\ref{lem:inc-det}
rules out a common nonzero center. Here $\br a b c=0$.
This identifies precisely why the converse cannot discard the
hypotheses that allow the triangle determinants to be canceled.


\section{Pascal: the universal identity and its constructive consequence}\label{sec:pascal}
Write $[ijk]=\det(P_i,P_j,P_k)$, with vectors as columns, and put
\[
 X=(P_1\times P_2)\times(P_4\times P_5),\quad
 Y=(P_2\times P_3)\times(P_5\times P_6),\quad
 Z=(P_3\times P_4)\times(P_6\times P_1).
\]
For $P=(x,y,z)$ let
$\nu(P)=(x^2,y^2,z^2,xy,xz,yz)$.  The universal identity is
\begin{equation}\label{eq:pascal-universal}
 \det(X,Y,Z)=\det\begin{pmatrix}
 \nu(P_1)\\\nu(P_2)\\\nu(P_3)\\
 \nu(P_4)\\\nu(P_5)\\\nu(P_6)
 \end{pmatrix}.
\end{equation}
This is an identity over $\mathbb Z$ in eighteen indeterminates.  Consequently
it remains valid after substitution into every commutative ring, including
rings of characteristic two.  The following proof never assumes that a
particular triple of input points is independent.

\subsection{The determinant relations used in the proof}
For five vectors $r,s,t,u,v$, the bracket relation is
\begin{equation}\label{eq:basic-plucker}
 [rst][ruv]-[rsu][rtv]+[rsv][rtu]=0.
\end{equation}
Here the letter bracket denotes the same three-column determinant.
To prove it without a nondegeneracy hypothesis, first expand the determinant
of the four-by-four matrix whose last three rows are the coordinates of
$r,s,t,u$ and whose first row repeats any one of those three rows.  Since
the matrix has a repeated row, its determinant is zero.  Expansion in the
first row, for each of its three possible choices, gives the vector identity
\[
 [stu]r-[rtu]s+[rsu]t-[rst]u=0.
\]
Taking its scalar product with $r\times v$ proves
\eqref{eq:basic-plucker}.  This proves the relation over every commutative
ring by finite determinant algebra.

The four particular instances needed below, with their signs displayed, are
\begin{align}
 [125][146]-[124][156]&=[126][145],\\
 [126][235]-[125][236]&=-[123][256],\\
 [134][236]-[234][136]&=[123][346],\\
 [245][156]-[145][256]&=-[125][456].
\end{align}
They result from taking respectively the common first vector in
\eqref{eq:basic-plucker} to be $P_1,P_2,P_3,P_5$ and reordering the other
arguments.  No division is involved.

\subsection{A two-term bracket expression}
The double-cross-product identity gives
\[
 X=[125]P_4-[124]P_5,\quad
 Y=[236]P_5-[235]P_6,\quad
 Z=[134]P_6-[346]P_1.
\]
On expanding the determinant, four of its eight terms vanish because they
have repeated columns.  Thus, writing $F=\det(X,Y,Z)$,
\begin{align*}
 F={}&[125][236][134][456]
       -[125][236][346][145]\\
    &+[125][235][346][146]
       -[124][235][346][156].
\end{align*}
The last two terms combine, by the first displayed bracket relation, to
$[235][346][126][145]$.  Combining that result with the second term and
using the second bracket relation yields
\begin{equation}\label{eq:pascal-two-term}
 F=[125][134][236][456]-[123][145][256][346].
\end{equation}

\subsection{Why the polynomial is alternating}
Interchanging $P_1$ and $P_2$ in \eqref{eq:pascal-two-term} and adding the
result to the original polynomial gives
\begin{align*}
 F+F_{(12)}={}&[125][456]\bigl([134][236]-[234][136]\bigr)\\
 &+[123][346]\bigl([245][156]-[145][256]\bigr)=0,
\end{align*}
by the third and fourth bracket relations.  Therefore $F$ changes sign
under the transposition $(12)$.

A cyclic shift of the six labels sends $(X,Y,Z)$ to $(Y,Z,-X)$.
Its effect on $F$ is consequently multiplication by $-1$.
Conjugating $(12)$ by powers of this cycle gives every adjacent
transposition; the two signs contributed by the conjugating powers cancel.
Thus $F$ changes sign under each adjacent transposition, and hence under any
permutation by that permutation's sign.  This is an exact polynomial
argument, with no appeal to a dense set or to a limiting configuration.

Each vertex occurs twice in every term of $F$.  Hence $F$ is homogeneous of
degree two in the three coordinates of each $P_i$ separately.  Each such
block has exactly the six possible monomials listed in $\nu(P_i)$.
Expand $F$ in products
\[
 \prod_{i=1}^6 \nu(P_i)_{k_i},\qquad k_i\in\{1,\ldots,6\}.
\]
If two indices $k_i,k_j$ are equal, swapping those two vertex blocks leaves
this monomial unchanged, whereas it negates its coefficient.  An integer
coefficient equal to its own negative is zero.  Thus every surviving tuple
$(k_1,\ldots,k_6)$ is a permutation of $(1,\ldots,6)$.
Antisymmetry further shows that its coefficient is the coefficient for
$(1,\ldots,6)$ multiplied by the sign of that permutation.  This is exactly
the determinant expansion.  Therefore $F=c\det(\nu(P_i))$ for one integer
$c$.

To determine $c$, specialize
\[
 P_1=(1,0,0),\quad P_2=(0,1,0),\quad P_3=(0,0,1),
\]
\[
 P_4=(u,v,w),\quad P_5=(x,y,z),\quad P_6=(p,q,r).
\]
Then
\[
 X=(uz-wx,vz-wy,0),\quad
 Y=(0,py-qx,pz-rx),\quad
 Z=(-qu,-qv,-rv).
\]
Their determinant is the following six-term polynomial:
\begin{align*}
 F={}&pquwyz-pqvwxz-pruvyz\\
     &+prvwxy+qruvxz-qruwxy.
\end{align*}
The coefficient of $(uv)(xz)(qr)$ is $+1$.
In the Veronese determinant the first three rows are the first three basis
rows, and its remaining minor has rows
$(uv,uw,vw)$, $(xy,xz,yz)$, $(pq,pr,qr)$; the same monomial is its diagonal
term, also of coefficient $+1$.  Hence $c=1$, proving
\eqref{eq:pascal-universal}.

The temporary cancellation of $2c=0$ in this proof occurs only in the
integer coefficient ring.  Once the resulting polynomial identity is
proved over $\mathbb Z$, specialization proves it in arbitrary commutative
rings.  In particular, this argument has not imposed a characteristic
restriction on the identity itself.

\subsection{From a conic equation to Pascal collinearity}
Let the six integer vectors satisfy
\[
 Q(P_i)=0\quad(1\leq i\leq6),\qquad
 Q(x,y,z)=q_1x^2+q_2y^2+q_3z^2+q_4xy+q_5xz+q_6yz,
\]
and suppose one specified coefficient $q_j$ is nonzero.  Put
$M_{ik}=\nu(P_i)_k$ and let $C_{ij}$ be its $(i,j)$ cofactor.
For each $j$ the completely explicit identity
\begin{equation}\label{eq:conic-certificate}
 q_j\det M=\sum_{i=1}^6 C_{ij}Q(P_i)
\end{equation}
holds.  Indeed, expand the right side in the coefficients $q_k$.
The coefficient of $q_k$ is $\sum_i C_{ij}M_{ik}$: this is $\det M$ when
$k=j$, by expansion down column $j$, and zero otherwise, by expansion of
the determinant obtained by replacing column $j$ with column $k$, which
has two identical columns.  The hypotheses make the right side zero.
Since $q_j\ne0$ and integer multiplication has no zero divisors,
$\det M=0$.  Equation~\eqref{eq:pascal-universal} gives
$\det(X,Y,Z)=0$.

This proves the finite-coordinate Pascal theorem even for a nonzero
singular conic.  Nonsingularity is not needed for the determinant
conclusion.  To read $X,Y,Z$ as actual projective intersection points,
the relevant adjacent vertices must define nonzero side lines and the
opposite side pairs must be distinct, so that the three double cross
products are nonzero.  The determinant statement remains valid when
these geometric objects degenerate, but the zero vector is not a
projective point. When $X,Y,Z$ are nonzero, Lemma~\ref{lem:axis}
constructs the actual nonzero common line $\axis(X,Y,Z)$. This completes
the proof of Theorem~\ref{thm:pas-summary}, including its witness.

Every rational configuration has integer homogeneous representatives:
multiply each point and the conic coefficients by the product of its
finitely many positive denominators.  This is a primitive-recursive
operation; no gcd, search for a root, or limiting process is required.
The same proof applies over a commutative integral domain, and over a
constructive field when the selected coefficient has an apartness witness.
For the literal first-order iPRA theorem, however, the integer-coded
formulation above is the one being asserted.

\subsection{An exact example with a point at infinity}
For the conic $xz-y^2=0$, take $P(t)=(t^2,t,1)$ for
$t=0,1,2,3,4,5$ in that order.  A chord through parameters $s,t$ has
representative $(1,-(s+t),st)$.  The three opposite-side intersections have
representatives
\[
 (2,2,1),\qquad(7,3,1),\qquad(5,1,0).
\]
They all satisfy the line equation $-x+5y-8z=0$.
The last point lies at infinity, which illustrates why homogeneous
coordinates remove any need for an exceptional ``parallel lines'' case.

\subsection{Verification artifact}
The dependency-free program \texttt{verify\_pascal.py} implements finite
integer-coefficient polynomials as dictionaries of sorted monomials.
It checks the universal determinant identity, the two-term bracket
formula, all four displayed Pluecker instances, and all six cofactor
certificates \eqref{eq:conic-certificate} with indeterminate conic
coefficients.  The Pascal polynomial has exactly $720$ nonzero monomials,
each with coefficient $+1$ or $-1$.  The program performs exact polynomial
normalization, not random sampling.  This is exact computational evidence; it is not a compiled Lean proof or
an emitted iPRA proof object.

\section{Why these proofs are in iPRA}\label{sec:internal}

\subsection{Explicit formulas and witnesses}
Every vector or coefficient is a fixed tuple of signed-integer pairs.
Addition, multiplication, equality, cross products, fixed-size
determinants, and $\axis$ are primitive recursive. Removing the
existential line from Pascal gives the following universally closed
formula:
\begin{align*}
 &\left(q\ne0\land\bigwedge_{i=1}^6Q(P_i)=0
             \land\nz(X)\land\nz(Y)\land\nz(Z)\right)\\
 &\quad\longrightarrow
 \left(\nz(h)\land h\cdot X=0\land h\cdot Y=0\land h\cdot Z=0\right),
 \qquad h=\axis(X,Y,Z).
\end{align*}
Here $h$ is an abbreviation for the displayed term, not an additional
existential quantifier. After expanding the abbreviations, the matrix
is quantifier free. The Desargues forward and reverse constructions
have the same form, with the supplied center or axis among the inputs.
The existential equivalence follows by ordinary intuitionistic
introduction and elimination of those witnesses.

All branch tests are decidable numerical predicates. An index $j<6$ is
selected by at most six tests, and a coordinate $i<3$ by at most three.
Alternatively the index can be supplied as a witness and its finitely
many possible values considered separately. No minimization over an
unbounded domain is used.

\subsection{Turning the algebra into actual derivations}
An equality of integer polynomials is an equality of finite coefficient
tables. Expanding both sides by distributivity and commutativity yields
the same signed monomial sum. This expansion is a finite derivation
from commutative-ring equations, not a semantic appeal to truth in the
standard natural numbers. Substitute signed-pair terms for the variables.
Section~\ref{sec:arith} supplies iPRA derivations of the ring equations,
so the finite derivation becomes an iPRA derivation.

This applies to the Desargues identity, the bracket identities, the
Pascal identity, and each of the six cofactor certificates. The
coefficient argument in Section~\ref{sec:pascal} explains Pascal's
identity to a human reader; direct expansion provides a fully finite
route for the formal proof as well. The supplied programs check this
coefficient equality with symbolic coordinates, rather than merely
checking numerical examples.

One may formalize a general compiler from polynomial identities to
object-theory derivations. That compiler is not logically necessary for
these fixed theorems: their fixed expansions can instead be generated
and replayed directly. In either case the formal target is
object-theory derivability, not merely acceptance by a Python program.

\subsection{Induction and logical resource audit}
\begin{center}
\begin{tabular}{p{0.30\linewidth}p{0.60\linewidth}}
\toprule
Component & Resources needed\\
\midrule
Natural and signed arithmetic & Defining equations and quantifier-free induction with parameters\\
Fixed determinant identities & Finite ring-equation derivations; no geometric induction\\
Selecting a nonzero entry & At most three or six decidable tests\\
Constructing a common line & The four branches of \eqref{eq:axis}\\
Integer cancellation & Decidable sign normalization and successor separation\\
Existential conclusions & Introduction of explicit primitive-recursive terms\\
PR real-name extension & First-order intuitionistic reasoning with quantified premises; no induction on those premises\\
\bottomrule
\end{tabular}
\end{center}

No compactness theorem, algebraic closure, root-finding principle, or
choice principle is required. In particular, this is not an argument
that all computable mathematical theorems must belong to iPRA: the actual
permitted derivation steps have been specified.

The coordinate bounds are also explicit. If the input point coordinates
have absolute value at most $H\ge1$, a cross product has coordinates at
most $2H^2$, and a double cross product at most $8H^4$. The first two
axis branches, applied to those double cross products, are bounded by
$128H^8$, and the other branches by $8H^4$. These coarse bounds are
primitive recursive and require no search for an output bound.


\section{A constructive real-coordinate schema}\label{sec:real}

The rational theorem already gives a literal iPRA interpretation of the
geometry. This extension explains precisely how the same proofs apply
to constructively given real coordinates, without adding a second-order
domain of arbitrary real numbers.

\subsection{Fixed primitive-recursive names and arithmetic predicates}
Fix, externally, finitely many primitive-recursive rational sequences
$a_i(n)$ with positive denominator functions. Positivity is either built
into the coding, as a denominator of the form $d_i(n)+1$, or included as
an explicit arithmetic premise. Their definitions may have numerical
parameters. Include as hypotheses, or prove in iPRA,
the regularity conditions
\begin{equation}\label{eq:regular}
 |a_i(m)-a_i(n)|\le2^{-m}+2^{-n}\qquad(m,n\in\N).
\end{equation}
Rational inequalities are interpreted by integer cross multiplication of
positive denominators, so these are first-order arithmetic formulas.
The instance at $m=0$ gives $|a_i(n)|\le|a_i(0)|+2$. For example,
$2$ plus the sum of the absolute initial numerators is an integer
bound $B\ge1$ for all the coordinates.

For each fixed integer polynomial $P$ in this finite tuple, write
$P_n=P(a_1(n),\ldots,a_d(n))$ and define
\begin{align}
 \Zero(P)&:\Longleftrightarrow
 \forall k\,\exists N\,\forall n\ge N\ (|P_n|\le2^{-k}),\label{eq:zero-real}\\
 \Apart(P)&:\Longleftrightarrow
 \exists r\,\exists N\,\forall n\ge N\ (|P_n|\ge2^{-r}).\label{eq:apart-real}
\end{align}
In the usual real interpretation, these mean that the polynomial value
is zero, or is positively apart from zero. More fundamentally here, they
are explicit arithmetic formulas. The statement is a schema for each
fixed finite list of PR definitions. It does not quantify over arbitrary
function codes or posit a total universal primitive-recursive evaluator.

\subsection{The estimates replacing exact cancellation}
If $P=\sum_\alpha c_\alpha x^\alpha$, set
\[
 M_P(B)=\sum_\alpha |c_\alpha|B^{|\alpha|}.
\]
Then $|P_n|\le M_P(B)$, term by term. All polynomials in use are fixed
and finite. iPRA proves the following rules from the regularity premises:
\begin{align}
 \Zero(P)\land\Zero(Q)&\Longrightarrow\Zero(P+Q),\label{eq:z-add}\\
 \Zero(P)&\Longrightarrow\Zero(SP),\label{eq:z-mul}\\
 \Apart(C)\land\Zero(CP)&\Longrightarrow\Zero(P).\label{eq:z-cancel}
\end{align}
For \eqref{eq:z-add}, given $k$, use both zero hypotheses at precision
$k+1$ and take the maximum of their supplied thresholds. The triangle
inequality gives $2^{-(k+1)}+2^{-(k+1)}=2^{-k}$.
For \eqref{eq:z-mul}, choose an explicit $t$ with $M_S(B)\le2^t$,
for example $t=M_S(B)+1$, and use the zero hypothesis at precision
$k+t$. The bound on $S$ gives the result.
For \eqref{eq:z-cancel}, open the witnesses $r,N_0$ of apartness.
Use $\Zero(CP)$ at precision $k+r$ and take the maximum of $N_0$ and
the resulting threshold. For every subsequent $n$,
\[
 2^{-r}|P_n|\le|C_nP_n|\le2^{-(k+r)},
\]
hence $|P_n|\le2^{-k}$. This is rational arithmetic with a positive dyadic
factor. Negation preserves $\Zero$, and fixed finite sums are obtained
by finitely many applications of \eqref{eq:z-add}. Apartness of two
factors gives apartness of their product by adding the exponents and
taking the maximum of the thresholds.

Each witness comes from a specified existential premise or a
primitive-recursive operation on finitely many witnesses. No sequence
of witnesses is chosen, and no induction is performed on $\Zero$ or
$\Apart$. Quantifier-free induction suffices for the elementary power
and rational-arithmetic bounds. These are consequently conditional
first-order iPRA derivations even though their hypotheses are quantified.

\subsection{Transferring the geometric proofs}
All polynomial identities hold exactly at every approximation index.
For Desargues, replace incidence through a center by
$\Zero(o\cdot j_i)$ and require $\Apart(o_j)$ for a designated center
coordinate. Equation~\eqref{eq:adj3}, finite sums, bounded multipliers,
and \eqref{eq:z-cancel} give $\Zero(\br{j_1}{j_2}{j_3})$.
The Desargues identity gives $\Zero(\br X Y Z)$.
To produce an actual nonzero line, require apartness of a designated
coordinate of $X\times Y$ and choose $h=X\times Y$. Its first two
incidences are polynomial identities; its third is the zero determinant.

For the converse, the supplied axis satisfies $\Zero(h\cdot X)$,
$\Zero(h\cdot Y)$, and $\Zero(h\cdot Z)$. Require apartness of both
triangle determinants and of a designated coordinate of this axis. The same argument
gives the zero joining-line determinant. Apartness of a designated
coordinate of $j_1\times j_2$ produces the nonzero center
$o=j_1\times j_2$.

For Pascal, assume $\Zero(Q(P_i))$ for all six vertices and
$\Apart(q_j)$ for a designated coefficient. The cofactor identity
\eqref{eq:conic-certificate} gives $\Zero(q_j\det M)$ by finite sums
and bounded multipliers. Cancellation yields $\Zero(\det M)$ and
therefore $\Zero(\br X Y Z)$. A designated apart nonzero coordinate
of $X\times Y$ again gives the Pascal line $h=X\times Y$.
The side and meet vectors have supplied apartness witnesses whenever
they are to denote projective objects.

The rational branch algorithm \eqref{eq:axis} is deliberately not used
on arbitrary real names: exact real equality is not being decided.
For this schema, apartness of a selected pair supplies the pivot.
These are explicit quantitative nondegeneracy assumptions for the
constructive real version; they are stronger than merely negating
equality to zero.

\subsection{Producing regular output names}
A raw polynomial sequence $P_n$ need not satisfy the unit modulus
\eqref{eq:regular}. The repair is explicit. Put
\[
 L_P(B)=\sum_{|\alpha|>0}|c_\alpha|\,|\alpha|B^{|\alpha|-1}.
\]
The telescoping product identity applied to each monomial proves
\[
 |P_m-P_n|\le L_P(B)(2^{-m}+2^{-n}).
\]
Choose $s$ with $L_P(B)\le2^s$, for example $s=L_P(B)+1$.
Then $\widetilde P(n)=P_{n+s}$ is regular in the sense of
\eqref{eq:regular}. For a vector use a common shift for its components.

The same arithmetic formulas $\Zero$ and $\Apart$ also apply to the
derived PR rational sequences below, including shifted differences;
equivalently, enlarge the fixed finite family of names to include their
specified shifts. To see that resampling does not change any incidence, first note
that $P_n-P_{n+s}$ tends to zero with the displayed explicit bound.
For several output components, their respective shifts give
coordinate differences tending to zero; telescoping products and
bounded multipliers then show that every fixed polynomial in the
original outputs and the corresponding polynomial in the resampled
outputs differ by a $\Zero$ sequence. Thus all proven zero equations
transfer. Apartness transfers as well: if $|P_n|\ge2^{-r}$ eventually,
use $|P_n-\widetilde P(n)|\le2^{-(r+1)}$ eventually to obtain
$|\widetilde P(n)|\ge2^{-(r+1)}$. Thresholds are combined by maxima.
Consequently the constructed line or center can be represented by
regular primitive-recursive real names without changing the conclusion.

This extension is a completed conditional proof, not a claim that every
classical real configuration has a primitive-recursive presentation or
that iPRA contains an unrestricted real-number sort.


\section{Lean formalization and the remaining work}\label{sec:lean}

\subsection{The repository interface actually inspected}
The repository was inspected at commit
\texttt{0bba41d70f061442357cae0f6de055434a33fb09}.
Its root toolchain is Lean \texttt{v4.32.0}, and its Lake configuration
pins Mathlib \texttt{v4.32.0}. Relevant inspected modules include
\texttt{PR.lean}, \texttt{Syntax.lean}, \texttt{TeqReflect.lean}, and
\texttt{NativeBiInterpretation.lean}, together with the ConciseIPRA
README. The \texttt{docs/drafts} listing was inspected; none of its listed
archives was needed as an additional premise for these proofs.

The existing bridge is the following theorem, shown with ASCII
renderings of the repository notation:
\begin{lstlisting}
theorem pra_iff_native
  {Gamma : List (Fm PTm n)} {phi : Fm PTm n} :
  (Gamma |- [PRAAx] phi) <->
  NPrf (Gamma.map NFm.ofPRA) (NFm.ofPRA phi)
\end{lstlisting}
Its namespace is \texttt{LeanProofs.ConciseIPRA}. The theorem and formula
translations are in \texttt{NativeBiInterpretation.lean}
\cite{repo-bridge}. Thus an object-iPRA derivation, once constructed,
already has a route into the concise native calculus. A new proof of
bi-interpretability is unnecessary.

\subsection{What the supplied Lean sources contain}
The archive has two concrete source candidates for the ambient algebra:
\texttt{GeometryCore.lean} and \texttt{PascalConic.lean}. They use
explicit three- and six-component structures, dot and cross products,
a six-term $3\times3$ determinant, and a 20-term Laplace expansion of
the six-row determinant. The main equalities have \texttt{ring} proof
scripts; this is Mathlib's commutative ring normalization tactic
\cite{mathlib-ring}.

The source includes the full universal Desargues identity, its
center-to-collinearity implication and converse determinant implication,
the explicit-axis incidence lemma for an apart pair, the Pascal bracket
and quadratic-evaluation determinant identities, and the general-conic
Pascal implication. The last uses six column-replacement identities:
replacing column $j$ of $M$ by $Mq$ gives determinant $q_j\det M$.
The six conic equations make that replaced column zero; cancellation
then gives $\det M=0$. This is the same adjugate argument as the paper,
written without depending on a high-level matrix or geometry API.

\textbf{Validation status.} No Lean or Lake executable was available
in the execution environment. These new sources have \emph{not} been
compiled or kernel-checked. They contain no proof holes or added axiom
declarations, but that is not a substitute for compilation. Elaboration,
API, or performance repairs may still be needed. In particular, no new
object-theory theorem of the form $[]\vdash[\mathrm{PRAAx}]\varphi$ is
claimed as an already checked artifact. The written mathematics and
executed exact polynomial certificates have a separate status.

\subsection{A concrete path to a complete Lean proof}
The following milestones specify the remaining obligations.
\begin{enumerate}
\item \textbf{Compile the ambient algebra.} In a Lean/Mathlib
\texttt{v4.32.0} environment, compile the supplied files and repair any
elaboration issues. Audit the resulting theorems with
\texttt{\#print axioms}, rejecting \texttt{sorryAx}. This checks the
algebra but does not yet construct an iPRA derivation.

\item \textbf{Implement the numerical codes in the object language.}
Using \texttt{PR.zero}, \texttt{PR.succ}, \texttt{PR.proj},
\texttt{PR.comp}, and \texttt{PR.prec}, define the symbols for addition,
multiplication, monus, equality tests, signed pairs, fixed tuples, and
the four-branch axis function. The inspected proof interface provides
\texttt{PRA.proj}, \texttt{PRA.comp}, \texttt{PRA.precZero},
\texttt{PRA.precSucc}, and \texttt{PRA.qfInd}.

\item \textbf{Prove the arithmetic library internally.}
Construct \texttt{Prf PRAAx} derivations of the lemmas in
Section~\ref{sec:arith}, including signed-ring laws, decidability, and
no-zero-divisors. These are derivations with arbitrary object variables;
correctness only in the standard model would not suffice.

\item \textbf{Replay finite algebra in the object theory.}
Either implement a derivation-producing normalizer for finite polynomial
expressions or generate the fixed distributivity, associativity,
commutativity, and cancellation chains used here. A reflective route
must prove its normalizer sound for \texttt{Prf PRAAx}; an ambient
\texttt{ring} proof does not automatically translate itself. The exact
coefficient checks fix the identities and signs to be replayed.

\item \textbf{Construct the geometric derivations.}
Encode the formulas of Section~\ref{sec:internal} as \texttt{Fm PTm}.
Replay the incidence lemmas, finite coefficient selection, cancellation,
and four axis branches. The existing intuitionistic deduction
constructors handle universal closure and explicit existential witnesses.
This gives the actual target objects
\texttt{Prf PRAAx [] desarguesSentence} and
\texttt{Prf PRAAx [] pascalSentence}, where the two sentences are to
be defined by the encoding just described.

\item \textbf{Translate to the native calculus.}
Apply \texttt{pra\_iff\_native.mp} to those object-theory derivations.
The output is a native \texttt{NPrf} derivation of each translated sentence.

\item \textbf{Optional real schema.}
For a fixed PR coordinate family, encode the regularity, zero, and
apartness formulas and prove the finite threshold calculations in
Section~\ref{sec:real}. Reuse the same polynomial certificates. This
extension is separate from the exact rational theorem.
\end{enumerate}

The substantial remaining implementation is the object-arithmetic
library and replay of finite ring-equation derivations. There is no
unresolved geometric lemma in this plan. It would nevertheless be
inaccurate to call the ambient Lean candidate a completed formal proof
in the deeply embedded iPRA.

\section{Executable evidence and reproducibility}\label{sec:evidence}

The archive contains three complementary checks of the formulas.
\begin{enumerate}
\item The dependency-free sparse-polynomial verifier checks Pascal's
identity, its two-product bracket form, the four bracket relations, and
all six conic cofactor certificates. Variables are indeterminates;
coefficients are exact integers. The Pascal polynomial has exactly
$720$ nonzero monomials, with coefficients $+1$ or $-1$.

\item A second symbolic verifier checks the complete Desargues identity
and all three coordinates of the adjugate identity
\eqref{eq:adj3}. The Desargues polynomial has $1{,}290$ nonzero monomials.
A supplementary checker verifies the six generic column-replacement
identities used by the Lean candidate.

\item The integer geometry program checks the explicit examples and
the common-kernel algorithm on the cube $\{-1,0,1\}^3$.
Of $18{,}954$ input triples with $u\ne0$, exactly $7{,}146$ have zero
determinant; all returned vectors are nonzero and orthogonal to all
three inputs. The four branches are used $5{,}040$, $1{,}872$, $216$,
and $18$ times. This finite check exercises the implementation; the
universal proof is Lemma~\ref{lem:axis}.
\end{enumerate}
The logs and polynomial fingerprints are included. These computations
verify exact formulas; they are not emitted iPRA proofs or substitutes
for Lean kernel checking. The source code uses the Python standard
library only.

Run the checks using the commands in \texttt{README.txt}. Rebuild the
paper from \texttt{paper.tex} with \texttt{pdflatex} twice, or run
\texttt{build.sh}. The archive includes the uncompiled Lean sources,
pinned build metadata, a detailed formalization roadmap, and source
URLs with hashes for the repository audit.

\section{Conclusions}
Both theorems have proofs whose arithmetic core is finite
commutative-ring algebra. Desargues is governed by the product of the
two triangle determinants and the determinant of the vertex-joining
lines. Pascal is governed by the determinant of the six quadratic
evaluation rows, with a nonzero conic coefficient canceled through an
explicit cofactor certificate. These calculations give complete iPRA
proofs for integer-coded rational configurations, including primitive
recursive witnesses for their geometric conclusions. The real-name
schema makes the extra apartness information explicit. The Lean plan
preserves the distinction between checking the algebra in Lean and
constructing derivations inside intuitionistic primitive recursive
arithmetic.


\begin{thebibliography}{9}
\bibitem{desargues}
\emph{Desargues's theorem}. Statement and geometric conventions supplied
in the task. \url{https://en.wikipedia.org/wiki/Desargues%27s_theorem}.
Accessed 15 November 2026. Used for the requested statement, not as a
premise of the proof.
\bibitem{pascal}
\emph{Pascal's theorem}. Statement supplied in the task.
\url{https://en.wikipedia.org/wiki/Pascal%27s_theorem}.
Accessed 15 November 2026. Used for the requested statement, not as a
premise of the proof.
\bibitem{repo-pr}
Vladimir Reshetnikov, \emph{ProveIt: ConciseIPRA, PR.lean}, commit
\texttt{0bba41d70f061442357cae0f6de055434a33fb09}.
\href{https://github.com/VladimirReshetnikov/ProveIt/blob/0bba41d70f061442357cae0f6de055434a33fb09/Logic/PrimitiveRecursiveArithmetic/ConciseIPRA/Lean/ConciseIPRA/PR.lean}{Pinned source: primitive-recursive symbols, terms, axioms, and quantifier-free induction}.
Inspected through the GitHub connector.
\bibitem{repo-bridge}
Vladimir Reshetnikov, \emph{ProveIt: ConciseIPRA,
NativeBiInterpretation.lean}, same commit.
\href{https://github.com/VladimirReshetnikov/ProveIt/blob/0bba41d70f061442357cae0f6de055434a33fb09/Logic/PrimitiveRecursiveArithmetic/ConciseIPRA/Lean/ConciseIPRA/NativeBiInterpretation.lean}{Pinned source: the faithful iPRA/native translation}.
Also inspected: \texttt{Syntax.lean}, \texttt{TeqReflect.lean}, the
ConciseIPRA README, root \texttt{lean-toolchain}, and root
\texttt{lakefile.toml} at the same commit.
\bibitem{mathlib-ring}
The Mathlib contributors, \emph{Mathlib.Tactic.Ring}, official
documentation and source.
\url{https://leanprover-community.github.io/mathlib4_docs/Mathlib/Tactic/Ring.html}.
The package targets the repository-pinned Mathlib version \texttt{v4.32.0}.
\end{thebibliography}
\end{document}

